Este capítulo expone la teoría de la cobertura con contratos de futuros, que fundamenta la hipótesis de la investigación.

\guia{El marco teórico presenta las teorías y la evidencia empírica que sustentan tu hipótesis. Su herramienta principal es la cita: cada idea tomada de otro autor se cita en formato APA 7.}

\section{Citas parentéticas y narrativas}

\guia{APA 7 usa el sistema autor-fecha. La cita parentética pone autor y año entre paréntesis al final de la idea; la cita narrativa integra el apellido en la oración y el año entre paréntesis. Con tres o más autores se escribe solo el primer apellido seguido de et al. desde la primera cita. Toda cita debe aparecer en la lista de referencias y viceversa. En LaTeX: \texttt{\textbackslash parencite\{clave\}} para la cita parentética y \texttt{\textbackslash textcite\{clave\}} para la narrativa; las claves se definen en \texttt{referencias.bib}.}

\textbf{Cita parentética:} los agentes con exposición en el mercado de contado toman posiciones en futuros para reducir la varianza de su riqueza \parencite{johnson1960}.

\textbf{Cita narrativa:} de acuerdo con \textcite{ederington1979}, la efectividad de una cobertura se mide como la proporción de la varianza de la posición que se elimina al cubrirla.

\textbf{Varios autores en una misma cita:} la teoría de la cobertura con futuros se basa en la reducción de la varianza de la posición de contado \parencite{johnson1960,ederington1979}.

\section{Citas textuales}

\guia{Una cita textual de menos de 40 palabras va entre comillas dentro del párrafo, con el número de página: (Apellido, año, p. 15). Si tiene 40 palabras o más, va en un bloque aparte, con sangría de 1.27 cm, sin comillas y con la cita al final. Prefiere parafrasear: en economía y finanzas las citas textuales se usan poco. En LaTeX: \texttt{\textbackslash parencite[p.~15]\{clave\}}; para el bloque, el entorno \texttt{quote}.}
